% Options for packages loaded elsewhere
\PassOptionsToPackage{unicode}{hyperref}
\PassOptionsToPackage{hyphens}{url}
\documentclass[
  a4paper,
]{ltjsarticle}
\usepackage{xcolor}
\usepackage[margin=25mm]{geometry}
\usepackage{amsmath,amssymb}
\setcounter{secnumdepth}{-\maxdimen} % remove section numbering
\usepackage{iftex}
\ifPDFTeX
  \usepackage[T1]{fontenc}
  \usepackage[utf8]{inputenc}
  \usepackage{textcomp} % provide euro and other symbols
\else % if luatex or xetex
  \usepackage{unicode-math} % this also loads fontspec
  \defaultfontfeatures{Scale=MatchLowercase}
  \defaultfontfeatures[\rmfamily]{Ligatures=TeX,Scale=1}
\fi
\usepackage{lmodern}
\ifPDFTeX\else
  % xetex/luatex font selection
\fi
% Use upquote if available, for straight quotes in verbatim environments
\IfFileExists{upquote.sty}{\usepackage{upquote}}{}
\IfFileExists{microtype.sty}{% use microtype if available
  \usepackage[]{microtype}
  \UseMicrotypeSet[protrusion]{basicmath} % disable protrusion for tt fonts
}{}
\makeatletter
\@ifundefined{KOMAClassName}{% if non-KOMA class
  \IfFileExists{parskip.sty}{%
    \usepackage{parskip}
  }{% else
    \setlength{\parindent}{0pt}
    \setlength{\parskip}{6pt plus 2pt minus 1pt}}
}{% if KOMA class
  \KOMAoptions{parskip=half}}
\makeatother
\usepackage{longtable,booktabs,array}
\newcounter{none} % for unnumbered tables
\usepackage{calc} % for calculating minipage widths
% Correct order of tables after \paragraph or \subparagraph
\usepackage{etoolbox}
\makeatletter
\patchcmd\longtable{\par}{\if@noskipsec\mbox{}\fi\par}{}{}
\makeatother
% Allow footnotes in longtable head/foot
\IfFileExists{footnotehyper.sty}{\usepackage{footnotehyper}}{\usepackage{footnote}}
\makesavenoteenv{longtable}
\setlength{\emergencystretch}{3em} % prevent overfull lines
\providecommand{\tightlist}{%
  \setlength{\itemsep}{0pt}\setlength{\parskip}{0pt}}
\usepackage{bookmark}
\IfFileExists{xurl.sty}{\usepackage{xurl}}{} % add URL line breaks if available
\urlstyle{same}
\hypersetup{
  pdftitle={無名複素波集合の因数分解としての創発構造 ――シンプレックス、局在、凝縮体および動力学の作業仮説と数値検証計画},
  pdfauthor={木原 範昭（Independent Researcher）　ORCID 0009-0004-6753-4020},
  hidelinks,
  pdfcreator={LaTeX via pandoc}}

\title{無名複素波集合の因数分解としての創発構造
――シンプレックス、局在、凝縮体および動力学の作業仮説と数値検証計画}
\author{木原 範昭（Independent Researcher）　ORCID 0009-0004-6753-4020}
\date{2026年9月2日　v1.0　Version DOI 10.5281/zenodo.22250744}

\begin{document}
\maketitle

\subsection{要旨}\label{ux8981ux65e8}

本稿は、空間、時間、粒子、頂点、辺および完全ネットワークを基礎構造として仮定せず、無名な複素波の多重集合だけから、それらに対応する構造が一意に読み出され得るかを問う作業仮説を定式化する。

従来、複数の要素間の関係を記述するとき、要素数を先に与え、その全二体関係数を完全グラフの辺数として数えることが多い。しかし、この記述は、識別可能で持続的な頂点、各関係の端点、完全な接続性および一関係一辺対応を暗黙に導入する。基礎にある波が無名であるなら、これらは公理ではなく、無名集合の因数分解に成功した後でのみ認められる再構成結果でなければならない。

本稿では、基礎状態を置換で同一視された複素波多重集合として定義する。その部分多重集合が、二乗ゼロ閉塞、有限回帰、整数比、共通位相、低ランク性、複素Gram整合性などを同時に満たして一意に因数分解された場合に限り、倍音、局在、凝縮体、シンプレックスまたは持続的対象として読む。完全ネットワークは入力ではなく、無名な関係量がシンプレックスへ一意に再構成された場合に、結果として判明する接続構造である。

さらに、動力学を、名前付き構成要素の軌跡ではなく、無名波集合の因数分解型が保存されつつ共通位相中心を変化させる過程として定義する。既知構造から生成したデータの由来情報を完全に消去する正解埋込み実験、構造を植え込まないブラインド対照実験、置換不変性試験、摂動試験および時間持続性試験を提案する。本稿は創発構造の存在または一意性を証明するものではない。無名性を破らずにそれらを検証するための反証可能な問題設定を固定するものである。

\textbf{キーワード}：無名複素波、因数分解、二乗ゼロ閉塞、有限回帰、シンプレックス、無ラベル距離幾何、局在、凝縮体、創発動力学、数値実験

\begin{center}\rule{0.5\linewidth}{0.5pt}\end{center}

\subsection{1. はじめに}\label{ux306fux3058ux3081ux306b}

物理モデルを構成するとき、粒子や節に番号を付け、その間の関係を配列または行列として記述する方法は自然である。要素数を
N とし、すべての二要素間に一つずつ関係を置けば、関係数は次式になる。

\[
M=\frac{N(N-1)}{2}.
\]

この数え上げ自体は正しい。しかし、その前には、少なくとも次の仮定が置かれている。

\begin{enumerate}
\def\labelenumi{\arabic{enumi}.}
\tightlist
\item
  N 個の節が先に存在する。
\item
  節どうしを区別できる。
\item
  節の同一性が変化の途中でも保存される。
\item
  すべての二節間に関係がある。
\item
  各関係は一組の端点を持つ。
\item
  三体以上の関係は二体関係から構成される。
\end{enumerate}

しかし、基礎にあるものが無名な複素波だけなら、波には最初から頂点番号も辺番号も付いていない。配列添字は計算上の記憶位置であって、物理的な名前ではない。それにもかかわらず、関係を最初から添字対で表せば、完全ネットワークと接続情報を基礎へ混入させることになる。

本稿は、この順序を逆転する。

\[
\text{完全ネットワーク}
\longrightarrow
\text{関係量}
\]

を出発点とせず、

\[
\text{無名な複素波集合}
\longrightarrow
\text{一意な因数分解}
\longrightarrow
\text{幾何・局在・凝縮体・動力学の読出し}
\]

を問う。

この変更により、完全ネットワークは存在論的前提ではなくなる。無名な関係量の一部が、ある頂点数を持つシンプレックスへ一意に再構成されたとき、初めてその部分集合が完全グラフの辺集合として読めたと判断する。

同じ考え方は倍音にも適用される。基底波とその子である倍音を先に置くのではない。無名波集合の中に、共通基本周期と整数比で因数分解できる部分集合が見つかったときに、結果として倍音系列だったと読む。局在、凝縮、親子関係および移動も同様に、無名集合の安定した因数分解結果として定義する。

\begin{center}\rule{0.5\linewidth}{0.5pt}\end{center}

\subsection{2.
本稿の立場と主張範囲}\label{ux672cux7a3fux306eux7acbux5834ux3068ux4e3bux5f35ux7bc4ux56f2}

本稿は証明論文でも数値結果報告でもない。次の三点を目的とする。

\begin{enumerate}
\def\labelenumi{\arabic{enumi}.}
\tightlist
\item
  無名性を保つ基礎状態を定義する。
\item
  創発構造を無名集合の因数分解として定義する。
\item
  仮説を循環なく検証し、失敗を判定できる実験手順を固定する。
\end{enumerate}

本稿では、空間、時間または粒子がすでに導出されたとは主張しない。特定の更新則が正しい物理法則であるとも主張しない。また、実在粒子、標準模型、重力または宇宙論との対応を主張しない。

本稿が提案する中心命題は次である。

\begin{quote}
無名な複素波多重集合の中に、強い整合性を持つ少数の因子が存在する場合、それらはシンプレックス、倍音、局在、凝縮体または持続的対象として、ラベルに依存せず一意に読み出される可能性がある。
\end{quote}

\begin{center}\rule{0.5\linewidth}{0.5pt}\end{center}

\subsection{3.
無名複素波状態}\label{ux7121ux540dux8907ux7d20ux6ce2ux72b6ux614b}

\subsubsection{3.1 基礎集合}\label{ux57faux790eux96c6ux5408}

計算上は M 個の複素波を並べて扱う。

\[
W=(w_1,w_2,\ldots,w_M),
\qquad
w_m\in\mathbb C.
\]

しかし、添字は物理的な名前ではない。任意の置換を同一状態とみなし、基礎状態を次の同値類として定義する。

\[
\mathcal W
=
[W]_{S_M}
=
\{(w_{\pi(1)},\ldots,w_{\pi(M)})\mid \pi\in S_M\}.
\]

したがって物理的入力は順序付きベクトルではなく、多重集合である。

\[
\mathcal W
=
\{w_1,\ldots,w_M\}/S_M.
\]

同じ値を持つ波が複数存在してもよいため、通常の集合ではなく多重集合として扱う。

\subsubsection{3.2
二乗ゼロ閉塞}\label{ux4e8cux4e57ux30bcux30edux9589ux585e}

本稿で採用する基礎拘束は、エルミートノルムではない複素二乗和の閉塞である。

\[
Q(\mathcal W)
=
\sum_{m=1}^{M}w_m^2
=0.
\]

この式は波の並び順に依存しない。また、任意の非零複素数による共通スケール変換に対して閉じている。

\[
Q(\lambda\mathcal W)
=
\lambda^2Q(\mathcal W)
=0.
\]

したがって、絶対スケールはこの条件だけからは定まらない。

\subsubsection{3.3 有限回帰}\label{ux6709ux9650ux56deux5e30}

有限回帰を、名前付きの各波が同じ場所へ戻る条件としてではなく、無名集合全体が置換同値を含めて元の状態へ戻る条件として定義する。

\[
U^q\mathcal W
\simeq
\mathcal W.
\]

ここで q
は有限回帰次数、記号の同値は少なくとも波の置換を含む。必要に応じて、共通位相、全体スケールまたは幾何学的合同変換も同値関係へ含める。

この定義では、一周期後に各配列要素が同じ添字へ戻る必要はない。集合として同じ状態が再構成されればよい。

\subsubsection{3.4
計算順序と創発時間の区別}\label{ux8a08ux7b97ux9806ux5e8fux3068ux5275ux767aux6642ux9593ux306eux533aux5225}

数値実験では更新を行うためのステップ番号が必要である。

\[
s=0,1,2,\ldots.
\]

しかし、この s
は計算機が処理した順序であり、物理的時間ではない。本稿ではこれを計算順序と呼ぶ。将来、時間として読める因子を構成する場合、それは無名波集合内の有限巡回または参照位相から別に再構成されなければならない。

\[
\tau_{\mathrm{read}}
=
F_{\mathrm{clock}}(\mathcal W_s,\mathcal W_{s+1},\ldots).
\]

\begin{center}\rule{0.5\linewidth}{0.5pt}\end{center}

\subsection{4. 無名因数分解}\label{ux7121ux540dux56e0ux6570ux5206ux89e3}

\subsubsection{4.1 因子の定義}\label{ux56e0ux5b50ux306eux5b9aux7fa9}

無名波状態から構造を読む操作を、因数分解写像として表す。

\[
\mathcal F:
\mathcal W
\longmapsto
\{F_1,F_2,\ldots,F_L\}.
\]

因子 F は、部分多重集合、潜在パラメータ、再構成規則および残差からなる。

\[
F
=
(C,\theta,R,\varepsilon),
\qquad
C\subseteq\mathcal W.
\]

ここで、C は因子を構成する無名部分多重集合、\(\theta\)
は共通周期、位相中心、整数次数、幾何次元などの潜在パラメータ、R
はそれらから C を再構成する規則、\(\varepsilon\) は再構成残差である。

厳密因子では次が成立する。

\[
\varepsilon(F)=0.
\]

数値実験では有限精度を考慮し、事前に固定した許容値以下を近似因子とする。

\[
\varepsilon(F)\leq\varepsilon_{\mathrm{tol}}.
\]

\subsubsection{4.2
同値な因数分解}\label{ux540cux5024ux306aux56e0ux6570ux5206ux89e3}

二つの因数分解が、入力波の置換、因子名の置換、全体位相、共通スケール、鏡像または採用した幾何の合同変換だけで異なる場合、それらを同一の物理解釈とみなす。

非同値な有効因数分解の個数を次で表す。

\[
\Omega_{\mathrm{fact}}(\mathcal W)
=
\#\left(
\frac{\{\text{有効な因数分解}\}}
{\text{許容対称性}}
\right).
\]

判定は次の三種類になる。

\[
\Omega_{\mathrm{fact}}=0
\quad\Longrightarrow\quad
\text{指定した構造として再構成不能},
\]

\[
\Omega_{\mathrm{fact}}=1
\quad\Longrightarrow\quad
\text{物理的同値を除いて一意に再構成},
\]

\[
\Omega_{\mathrm{fact}}>1
\quad\Longrightarrow\quad
\text{複数の非同値な構造解釈が存在}.
\]

植え込んだ正解が一つ見つかっただけでは存在を示すにすぎない。一意な創発を主張するには、非同値な別解がないことを確認しなければならない。

\begin{center}\rule{0.5\linewidth}{0.5pt}\end{center}

\subsection{5.
完全ネットワークを入力しないシンプレックス再構成}\label{ux5b8cux5168ux30cdux30c3ux30c8ux30efux30fcux30afux3092ux5165ux529bux3057ux306aux3044ux30b7ux30f3ux30d7ux30ecux30c3ux30afux30b9ux518dux69cbux6210}

\subsubsection{5.1
無名な複素二乗長集合}\label{ux7121ux540dux306aux8907ux7d20ux4e8cux4e57ux9577ux96c6ux5408}

ある部分多重集合を、複素二乗長候補として取り出す。

\[
D
=
\{d_1^2,d_2^2,\ldots,d_r^2\}.
\]

この段階では、頂点も辺も存在せず、どの長さがどの二点間に対応するかも与えない。

候補頂点数 n が存在するなら、必要な長さ数は次を満たす。

\[
r
=
\frac{n(n-1)}{2}.
\]

ここでこの式は基礎公理ではない。無名部分集合が n
頂点シンプレックスの全辺候補になり得るかを判定するための必要条件である。

\subsubsection{5.2
一時的な辺割当て}\label{ux4e00ux6642ux7684ux306aux8fbaux5272ux5f53ux3066}

再構成器は、無名長集合から候補辺への全ての本質的に異なる割当てを調べる。

\[
f:
D
\longrightarrow
\bigl\{\{i,j\}\mid 0\leq i<j<n\bigr\}.
\]

この写像は基礎データではなく、逆問題の候補変数である。

基準頂点を 0 として、割当て f から複素Gram行列を構成する。

\[
G_{ab}
=
\frac{1}{2}
\left(
d_{0a}^2+d_{0b}^2-d_{ab}^2
\right),
\qquad
1\leq a,b<n.
\]

要求する埋込み次元、rank、非退化性、閉塞条件および追加の対称条件を満たす割当てだけを有効とする。複素二乗長を扱うため、実ユークリッド距離に対する正値性や三角不等式を無条件には要求しない。どの合同変換を同一視するかは、採用する複素双線形形式とともに実験仕様へ明記する。

\subsubsection{5.3
完全ネットワークは再構成証明書である}\label{ux5b8cux5168ux30cdux30c3ux30c8ux30efux30fcux30afux306fux518dux69cbux6210ux8a3cux660eux66f8ux3067ux3042ux308b}

頂点置換や合同変換を除いて有効割当てが一つだけ残った場合に限り、D は n
頂点シンプレックスを一意に定めたと判定する。

\[
D
\xrightarrow{\text{一意再構成}}
\text{n 頂点シンプレックス}.
\]

このとき初めて、D の要素が完全グラフの全辺に対応していたことが分かる。

\[
K_n
\quad\text{は入力ではなく再構成結果である}.
\]

実数距離の無ラベル集合から点配置を再構成する問題は、無ラベル距離幾何および無ラベル大域剛性として既に研究されている。本稿はその問題意識を共有するが、入力を正の実距離に限定せず、複素波、二乗ゼロ閉塞、有限回帰、倍音因子および因子の持続性を同一の読出し枠組みで扱う点を作業仮説として追加する。

\begin{center}\rule{0.5\linewidth}{0.5pt}\end{center}

\subsection{6.
倍音、局在および凝縮の因数分解}\label{ux500dux97f3ux5c40ux5728ux304aux3088ux3073ux51ddux7e2eux306eux56e0ux6570ux5206ux89e3}

\subsubsection{6.1
倍音は先験的な親子関係ではない}\label{ux500dux97f3ux306fux5148ux9a13ux7684ux306aux89aaux5b50ux95a2ux4fc2ux3067ux306fux306aux3044}

基底波を先に置き、その整数倍を子波として生成するのではない。無名波集合の位相履歴または巡回特性から、ある部分集合
C に対して共通基本量と整数列が存在するかを探索する。

\[
\omega_m
\simeq
k_m\omega_0,
\qquad
k_m\in\mathbb Z,
\qquad
w_m\in C.
\]

基本量に対応する波が集合内に明示的に存在する必要はない。基本量は複数の波に共通する潜在因子として逆算され得る。

位相についても、次の低次元表現を探索する。

\[
\phi_m(s)
\simeq
k_m\phi_0(s)+\delta_{k_m}
\pmod{2\pi}.
\]

この表現が小さな残差で成立した場合、C を倍音因子として読む。

\subsubsection{6.2
局在は共通位相の読出しである}\label{ux5c40ux5728ux306fux5171ux901aux4f4dux76f8ux306eux8aadux51faux3057ux3067ux3042ux308b}

倍音因子を合成して得られる読出し核を次のように書く。

\[
L_C(\phi)
=
\left|
\sum_{m\in C}
A_m
\exp\left[
i k_m(\phi-\phi_0)
+i\delta_{k_m}
\right]
\right|^2.
\]

L が位相 \(\phi_0\)
付近で鋭いピークを持つとき、その因子は局在状態として読める。ただし、これは基礎波が背景空間内の一点に存在するという意味ではない。無名波集合の共通位相因子が、ある読出しに対して鋭い一致位置を持つという意味である。

最高倍音次数を K
とすると、等振幅で規則的な位相を持つ理想的な系列では、位相分解能は概ね次の次数で改善する。

\[
\Delta\phi
\sim
\frac{1}{K}.
\]

多数の倍音が存在しても、それらが一つの共通基本位相と固定された整数列から生成されるなら、独立自由度が倍音数と同じだけ増えたとは限らない。多数の成分が少数の潜在因子で記述できること自体が、凝縮または低ランク性の指標となる。

\subsubsection{6.3 凝縮体}\label{ux51ddux7e2eux4f53}

本稿では凝縮体を、名前付き要素の固定集合ではなく、次を満たす持続的な低次元因子として定義する。

\begin{enumerate}
\def\labelenumi{\arabic{enumi}.}
\tightlist
\item
  多数の無名波を少数の潜在パラメータで再構成できる。
\item
  部分二乗閉塞または指定した保存残差が小さい。
\item
  共通位相中心と有限の位相幅を持つ。
\item
  許容対称性を除いた因子署名が計算順序をまたいで保存される。
\end{enumerate}

凝縮体の構成波が更新の途中で入れ替わってもよい。保存対象は要素名ではなく因数分解型である。

\begin{center}\rule{0.5\linewidth}{0.5pt}\end{center}

\subsection{7.
対称性と一意性}\label{ux5bfeux79f0ux6027ux3068ux4e00ux610fux6027}

\subsubsection{7.1
強い対称性の必要性}\label{ux5f37ux3044ux5bfeux79f0ux6027ux306eux5fc5ux8981ux6027}

二乗ゼロ閉塞、有限回帰、整数比、低ランク性および幾何整合性を同時に満たす集合は、一般の無作為な複素波集合の中では特殊であると予想される。強い対称性は、無名な多数成分を少数因子へ圧縮可能にする候補条件である。

一方、完全な対称性は、どの頂点、方向または位相中心も区別できない縮退を生む。したがって、必要なのは単純に最大の対称性ではない。

\[
\text{因数分解を成立させる強い対称性}
+
\text{読出しを選択する小さな対称性の破れ}
\]

という二層構造を作業仮説とする。

\subsubsection{7.2
対称解の同一視}\label{ux5bfeux79f0ux89e3ux306eux540cux4e00ux8996}

正単体の頂点名を変えて得られる多数の解は、異なる物理構造ではない。したがって、一意性はラベル付き解の個数ではなく、許容対称性で商を取った同値類の個数で判定する。

しかし、対称変換で結べない異なるシンプレックス、異なる倍音基本量または異なる凝縮分割が残る場合、それは真の非一意性である。

\subsubsection{7.3 摂動領域}\label{ux6442ux52d5ux9818ux57df}

完全対称状態を基準に、無名性を保つ摂動を加える。

\[
w_m(\epsilon)
=
w_m^{(0)}+\epsilon\eta_m.
\]

次の三領域が存在する可能性を検証する。

\begin{enumerate}
\def\labelenumi{\arabic{enumi}.}
\tightlist
\item
  完全対称点：形状は明瞭だが方向または中心が縮退する。
\item
  小摂動領域：形状を保存しつつ、位置または方向が選択される。
\item
  大摂動領域：有効因子が失われ、再構成が非一意または不能になる。
\end{enumerate}

臨界摂動量の存在自体が作業仮説であり、実験で否定され得る。

\begin{center}\rule{0.5\linewidth}{0.5pt}\end{center}

\subsection{8.
因子型の保存としての動力学}\label{ux56e0ux5b50ux578bux306eux4fddux5b58ux3068ux3057ux3066ux306eux52d5ux529bux5b66}

\subsubsection{8.1
微視的更新と観測動力学}\label{ux5faeux8996ux7684ux66f4ux65b0ux3068ux89b3ux6e2cux52d5ux529bux5b66}

微視的更新は、無名状態の列を作る写像として表す。

\[
\mathcal W_{s+1}
=
\Phi(\mathcal W_s).
\]

更新写像は入力順序へ依存してはならない。

\[
\Phi(\pi\mathcal W)
\simeq
\Phi(\mathcal W).
\]

各計算順序で因数分解を行う。

\[
\mathcal F(\mathcal W_s)
=
\{F_1(s),F_2(s),\ldots\}.
\]

観測される動力学は、個々の波の軌跡ではなく、因子同値類の変化として定義する。

\[
[F_\alpha(s)]
\longrightarrow
[F_\beta(s+1)].
\]

\subsubsection{8.2 因子署名}\label{ux56e0ux5b50ux7f72ux540d}

因子の同一性を判定するため、要素名に依存しない署名を定義する。

\[
\Sigma(F)
=
\left(
q,
\{k_m\},
\text{振幅比},
\text{閉塞残差},
\text{Gramスペクトル},
\text{rank},
\text{局在幅},
\ldots
\right).
\]

具体的な署名成分は対象因子ごとに事前登録する。更新前後で署名が一致し、共通位相中心だけが変化する場合、その因子を保存された凝縮体として読む。

\subsubsection{8.3 移動}\label{ux79fbux52d5}

因子形状を \(C_0\)、共通位相移動を T
とすれば、移動する因子は次で表される。

\[
F(s)
\simeq
T_{\phi(s)}F^{(0)}.
\]

更新後も次が成立するなら、塊は保存されている。

\[
\Sigma(F(s+1))
\simeq
\Sigma(F(s)).
\]

一方、中心位相は変化してよい。

\[
\phi(s+1)
\neq
\phi(s).
\]

したがって移動とは、同じ名前の材料が位置を変えることではなく、同じ因数分解型が共通位相中心を変えながら持続することである。

\subsubsection{8.4
融合、分裂および相互作用}\label{ux878dux5408ux5206ux88c2ux304aux3088ux3073ux76f8ux4e92ux4f5cux7528}

因数分解の組替えを次のように読む。

\[
\{F_A,F_B\}
\longrightarrow
\{F_C\}
\qquad
\text{融合候補},
\]

\[
\{F_A\}
\longrightarrow
\{F_B,F_C\}
\qquad
\text{分裂候補},
\]

\[
\{F_A,F_B\}
\longrightarrow
\{F'_A,F'_B\}
\qquad
\text{散乱または交換候補}.
\]

これらを物理的相互作用と呼ぶためには、因子数の変化とは独立に保存される全体不変量を特定しなければならない。本稿では、その具体形を仮定しない。

\begin{center}\rule{0.5\linewidth}{0.5pt}\end{center}

\subsection{9. 作業仮説}\label{ux4f5cux696dux4eeeux8aac}

本稿の主張を、将来の実験で個別に否定できる形へ分ける。

\subsubsection{仮説
H1：無名シンプレックス再構成}\label{ux4eeeux8aac-h1ux7121ux540dux30b7ux30f3ux30d7ux30ecux30c3ux30afux30b9ux518dux69cbux6210}

強い整合性を持つ一部の無名複素二乗量集合は、頂点置換、合同変換および許容ゲージを除いて、一つの複素シンプレックスへ一意に再構成される。

\subsubsection{仮説
H2：完全ネットワークの事後的出現}\label{ux4eeeux8aac-h2ux5b8cux5168ux30cdux30c3ux30c8ux30efux30fcux30afux306eux4e8bux5f8cux7684ux51faux73fe}

一意なシンプレックス再構成が成立した部分集合では、その要素数と接続割当てから、結果として完全グラフ構造が判明する。完全ネットワークを入力する必要はない。

\subsubsection{仮説
H3：倍音因子による局在読出し}\label{ux4eeeux8aac-h3ux500dux97f3ux56e0ux5b50ux306bux3088ux308bux5c40ux5728ux8aadux51faux3057}

整数比、共通基本位相および規則的振幅を共有する無名波部分集合は、合成読出しにおいて有限幅の位相ピークを作る。このピーク中心と幅が、局在位置と局在幅として読まれる。

\subsubsection{仮説
H4：低ランク因子としての凝縮体}\label{ux4eeeux8aac-h4ux4f4eux30e9ux30f3ux30afux56e0ux5b50ux3068ux3057ux3066ux306eux51ddux7e2eux4f53}

多数の無名波が少数の潜在因子で再構成可能であり、閉塞および位相幅が安定する場合、その因子は凝縮体として持続する。

\subsubsection{仮説
H5：因子型の保存としての移動}\label{ux4eeeux8aac-h5ux56e0ux5b50ux578bux306eux4fddux5b58ux3068ux3057ux3066ux306eux79fbux52d5}

凝縮因子の署名が保存され、共通位相中心だけが連続的に変化する場合、その変化は移動として読まれる。

\subsubsection{仮説
H6：創発的な空間・時間・対象}\label{ux4eeeux8aac-h6ux5275ux767aux7684ux306aux7a7aux9593ux6642ux9593ux5bfeux8c61}

シンプレックス因子、内部巡回因子および持続凝縮因子が同じ無名波系列から独立に再構成され、互いに整合する場合、それらはそれぞれ空間、時間および対象として読むことができる。

\subsubsection{仮説
H7：対称性と小さな破れ}\label{ux4eeeux8aac-h7ux5bfeux79f0ux6027ux3068ux5c0fux3055ux306aux7834ux308c}

一意な因数分解を可能にする強い対称性と、位置・方向を選択する小さな対称性の破れが両立する有限領域が存在する。

\begin{center}\rule{0.5\linewidth}{0.5pt}\end{center}

\subsection{10.
数値検証計画}\label{ux6570ux5024ux691cux8a3cux8a08ux753b}

\subsubsection{10.1
生成器と復元器の完全分離}\label{ux751fux6210ux5668ux3068ux5fa9ux5143ux5668ux306eux5b8cux5168ux5206ux96e2}

実験は二つの独立プログラムで構成する。

生成器は既知構造から複素波集合を作る。

\[
S_{\mathrm{true}}
\xrightarrow{\mathrm{generator}}
\mathcal W.
\]

生成後、次を全て削除する。

\begin{itemize}
\tightlist
\item
  頂点番号
\item
  辺端点
\item
  波の生成順序
\item
  倍音次数
\item
  基本波ラベル
\item
  凝縮体ラベル
\item
  親子ラベル
\item
  計算順序をまたぐ個体ID
\item
  生成に使用した座標
\end{itemize}

さらに入力順を無作為化する。

復元器へ渡すのは、許可された複素値、位相履歴および全体拘束だけである。復元器は生成器の内部データ、正解ラベルまたは接続表へアクセスしてはならない。

\subsubsection{10.2 実験
E0：置換不変性試験}\label{ux5b9fux9a13-e0ux7f6eux63dbux4e0dux5909ux6027ux8a66ux9a13}

同一の無名集合を複数のランダム置換で入力する。

\[
\mathcal W^{(r)}
=
\{w_{\pi_r(1)},\ldots,w_{\pi_r(M)}\}.
\]

すべての入力で、許容対称性を除いて同じ因数分解結果が得られることを要求する。非同値な結果が入力順序へ依存する場合、実装に隠れたラベル依存がある。

\subsubsection{10.3 実験
E1：正解埋込みシンプレックス}\label{ux5b9fux9a13-e1ux6b63ux89e3ux57cbux8fbcux307fux30b7ux30f3ux30d7ux30ecux30c3ux30afux30b9}

既知の実または複素シンプレックスから全二点関係量を生成し、対応表を消去する。復元器が次を満たすか調べる。

\begin{enumerate}
\def\labelenumi{\arabic{enumi}.}
\tightlist
\item
  元のシンプレックスと合同な解を回収する。
\item
  非同値な別解が存在しない。
\item
  頂点数を入力しない条件でも候補頂点数を回収できる。
\item
  完全ネットワークが出力結果として得られる。
\end{enumerate}

\subsubsection{10.4 実験
E2：正解埋込み倍音・局在}\label{ux5b9fux9a13-e2ux6b63ux89e3ux57cbux8fbcux307fux500dux97f3ux5c40ux5728}

既知の基本位相と整数次数から波集合を生成し、それらを無名化する。復元器が、基本量、整数次数集合、位相中心および局在幅を回収できるか調べる。

基本量に対応する波を入力集合から除いた場合も試験し、潜在基本因子を逆算できるかを確認する。

\subsubsection{10.5 実験
E3：複合因子}\label{ux5b9fux9a13-e3ux8907ux5408ux56e0ux5b50}

同一の無名波集合に、シンプレックス因子、倍音因子および閉塞因子を同時に埋め込む。個別に求めた因子が、同じ部分集合または整合する重なりを持つかを検証する。

この実験は、異なる読出しが同じ創発構造を指すかを調べる。

\subsubsection{10.6 実験
E4：持続・移動ロック}\label{ux5b9fux9a13-e4ux6301ux7d9aux79fbux52d5ux30edux30c3ux30af}

因子形状を保存し、共通位相中心だけを変化させた無名集合列を生成する。各計算順序で波順序と構成要素を交換し、名前付き追跡を不可能にする。

復元器が、構成波の同一性を使わずに同じ因子署名と位相中心の変化を回収できるかを検証する。

\subsubsection{10.7 実験
E5：ブラインド対照}\label{ux5b9fux9a13-e5ux30d6ux30e9ux30a4ux30f3ux30c9ux5bfeux7167}

構造を植え込まず、二乗ゼロ閉塞、有限回帰または同じ周辺分布だけを満たす無名集合を生成する。正解埋込み群と比較して、有効因子、一意因子および偽陽性の頻度を測る。

\[
P_{\mathrm{unique}}
=
P\left(
\Omega_{\mathrm{fact}}=1
\mid
\text{生成条件}
\right).
\]

正解埋込み群と対照群で同じ頻度の構造が検出されるなら、復元条件が弱すぎる。

\subsubsection{10.8 実験
E6：摂動と臨界性}\label{ux5b9fux9a13-e6ux6442ux52d5ux3068ux81e8ux754cux6027}

正解埋込み集合へ制御された摂動を加え、次を測る。

\begin{itemize}
\tightlist
\item
  因数分解数
\item
  正解回収率
\item
  再構成残差
\item
  Gram rank とスペクトル
\item
  局在ピーク幅
\item
  因子署名の寿命
\item
  偽陽性率
\end{itemize}

完全対称点から因子崩壊までの領域を連続的に調べる。

\begin{center}\rule{0.5\linewidth}{0.5pt}\end{center}

\subsection{11. 成功条件}\label{ux6210ux529fux6761ux4ef6}

本研究プログラムの最小成功条件を次のように定める。

\begin{enumerate}
\def\labelenumi{\arabic{enumi}.}
\tightlist
\item
  入力順序を変えても同じ物理同値類が得られる。
\item
  生成時の名前、座標、接続情報を使わず正解構造を回収できる。
\item
  正解と対称同値でない別解が存在しない。
\item
  複数の独立した読出しが同じ因子を指す。
\item
  小さな摂動に対して因子が有限の安定領域を持つ。
\item
  構成波を追跡せず、因子署名の保存だけで持続対象を同定できる。
\item
  計算ステップそのものを物理時間とせず、内部巡回から時間候補を再構成できる。
\item
  因子状態から次の因子状態について検証可能な予測が得られる。
\end{enumerate}

空間、時間および対象についてそれぞれ別の読出しを作るだけでは不十分である。それらが同一の無名波集合上で整合し、互いの予測を制約することを要求する。

\begin{center}\rule{0.5\linewidth}{0.5pt}\end{center}

\subsection{12. 反証条件}\label{ux53cdux8a3cux6761ux4ef6}

次のいずれかが成立した場合、対応する仮説は支持されない。

\begin{enumerate}
\def\labelenumi{\arabic{enumi}.}
\tightlist
\item
  入力配列の順序によって非同値な構造が出る。
\item
  正解を回収できても、同程度に有効な非同値解が多数残る。
\item
  正解ラベルまたは接続表を復元器へ渡さなければ構造を回収できない。
\item
  ランダム対照群でも正解埋込み群と同程度に構造が検出される。
\item
  任意に小さな摂動で因子が完全に失われる。
\item
  局在読出しが生成時に使用した外部座標へ依存する。
\item
  移動の判定が名前付き構成波の追跡へ依存する。
\item
  創発時間が計算ステップの別名にすぎない。
\item
  シンプレックス読出し、倍音局在読出しおよび持続因子読出しが互いに矛盾する。
\item
  探索範囲を拡張すると一意性が失われ、以前の一意解が探索不足による見かけだったと判明する。
\end{enumerate}

失敗結果も重要である。どの条件で因数分解が不可能、非一意または不安定になるかは、創発構造の成立範囲を定める。

\begin{center}\rule{0.5\linewidth}{0.5pt}\end{center}

\subsection{13.
計算量と探索戦略}\label{ux8a08ux7b97ux91cfux3068ux63a2ux7d22ux6226ux7565}

無名な r 個の候補長を、n
頂点シンプレックスの辺へ割り当てる単純探索は急速に増大する。全ての値が異なる一般の場合、頂点置換を除く候補数の概算は次である。

\[
\frac{r!}{n!},
\qquad
r=\frac{n(n-1)}{2}.
\]

したがって、大規模系では全順列探索ではなく、次の枝刈りが必要になる。

\begin{itemize}
\tightlist
\item
  部分Gram整合性
\item
  rank条件
\item
  二乗ゼロ閉塞する部分集合
\item
  有限巡回次数
\item
  整数比候補
\item
  重複値と自己同型
\item
  スペクトル不変量
\item
  低ランク分解
\item
  部分シンプレックスの貼合せ整合性
\end{itemize}

正解埋込み実験では生成器が正解を知っているが、復元器がその正解を枝刈りへ使用してはならない。正解は復元終了後の評価にだけ使う。

探索が大規模化した場合、全解列挙の代わりに一意性証明書を構成する方法も検討する。ただし、局所最適解を一つ発見しただけで一意性を主張してはならない。

\begin{center}\rule{0.5\linewidth}{0.5pt}\end{center}

\subsection{14.
既存研究との最小限の関係}\label{ux65e2ux5b58ux7814ux7a76ux3068ux306eux6700ux5c0fux9650ux306eux95a2ux4fc2}

実数点配置が無ラベルな距離多重集合から再構成できるかという問題には既存研究がある。Boutin
と Kemper
は、距離または面積の分布から点配置を再構成する問題を扱い、多くの一般的位置配置が再構成可能である一方、反例も存在することを示した。Gortler、Theran、Thurston
は、無ラベル距離と大域剛性の一般的位置での関係を研究した。

本稿はこれらの結果を複素波系へそのまま適用したとは主張しない。実正距離に対する正値性を持たない複素二乗量、二乗ゼロ閉塞、有限回帰、倍音因子および時間をまたぐ因子持続性を同時に扱うには、別の定式化と数値検証が必要である。

\begin{center}\rule{0.5\linewidth}{0.5pt}\end{center}

\subsection{15.
限界と今後の課題}\label{ux9650ux754cux3068ux4ecaux5f8cux306eux8ab2ux984c}

本稿の段階では、次が未解決である。

\begin{enumerate}
\def\labelenumi{\arabic{enumi}.}
\tightlist
\item
  無名性として置換だけを商に取るべきか、より広い基底変換も同一視すべきか。
\item
  複素二乗長の合同群と、一意再構成の正確な同値関係。
\item
  因子が重なり合う場合の全体構造が、分割、被覆またはハイパーグラフのどれになるか。
\item
  有効因子残差の自然な尺度と許容値。
\item
  強い対称性が一意因数分解を増やす条件と、逆に縮退を増やす条件。
\item
  内部巡回から創発時間を構成する具体的手順。
\item
  因子間相互作用を生む、無名性を保った微視的更新則。
\item
  因子数が変化しても保存される全体不変量。
\item
  大規模探索を可能にする一意性証明アルゴリズム。
\end{enumerate}

これらは本稿の欠落を隠すものではなく、作業仮説を次の実験へ接続するための明示的課題である。

\begin{center}\rule{0.5\linewidth}{0.5pt}\end{center}

\subsection{16. 結論}\label{ux7d50ux8ad6}

本稿は、完全ネットワーク、倍音親子、局在体または粒子を基礎へ置かず、無名な複素波多重集合だけからそれらを読み出す研究枠組みを提案した。

\[
\mathcal W
=
\{w_1,\ldots,w_M\}/S_M
\]

を基礎状態とし、そこから有効な部分因子を探索する。

\[
\mathcal W
\xrightarrow{\text{無名因数分解}}
\left\{
\begin{array}{l}
\text{シンプレックス因子}\\
\text{倍音因子}\\
\text{局在因子}\\
\text{凝縮因子}\\
\text{持続因子}
\end{array}
\right\}.
\]

シンプレックス因子が一意に再構成された場合、完全ネットワークは入力ではなく、その再構成結果として判明する。倍音関係も親子構造を先に置かず、共通基本量と整数比を共有する部分集合の因数分解結果として現れる。局在は共通位相因子の鋭い読出しであり、凝縮体は多数の無名波を少数の潜在量で記述できる低ランク因子である。

動力学は、名前付き波の軌跡ではなく、因子署名を保存しながら共通位相中心を変える因子同値類の変化として定義される。

\[
[F(s)]
\longrightarrow
[F(s+1)].
\]

この枠組みの成否は、正解埋込みデータを無名化したブラインド復元、構造を植え込まない対照実験、置換不変性試験、摂動試験および持続性試験によって判定できる。

本稿の中心的な提案は、構造を仮定して波へ割り当てることではない。

\[
\boxed{
\begin{array}{l}
\text{無名な波集合が、どの条件で一つの構造として}\\
\text{一意に因数分解されるかを問う}
\end{array}
}
\]

ことである。空間、時間および粒子として読める因子が、生成時の名前や座標を使わず、同一の無名波集合から互いに整合して再構成されるなら、本研究プログラムの最小目標は達成される。

\begin{center}\rule{0.5\linewidth}{0.5pt}\end{center}

\subsection{参考文献}\label{ux53c2ux8003ux6587ux732e}

\begin{enumerate}
\def\labelenumi{\arabic{enumi}.}
\item
  Mireille Boutin and Gregor Kemper, ``On Reconstructing n-Point
  Configurations from the Distribution of Distances or Areas,''
  \emph{Advances in Applied Mathematics}, Vol. 32, No.~4, pp.~709--735,
  2004. \url{https://doi.org/10.1016/S0196-8858(03)00101-5}
\item
  Steven J. Gortler, Louis Theran, and Dylan P. Thurston, ``Generic
  Unlabeled Global Rigidity,'' \emph{Forum of Mathematics, Sigma}, Vol.
  7, e21, 2019. \url{https://doi.org/10.1017/fms.2019.16}
\end{enumerate}

\begin{center}\rule{0.5\linewidth}{0.5pt}\end{center}

\subsection{付録A：記号一覧}\label{ux4ed8ux9332aux8a18ux53f7ux4e00ux89a7}

{\def\LTcaptype{none} % do not increment counter
\begin{longtable}[]{@{}ll@{}}
\toprule\noalign{}
記号 & 意味 \\
\midrule\noalign{}
\endhead
\bottomrule\noalign{}
\endlastfoot
M & 基礎入力に含まれる無名複素波の総数 \\
W & 計算上の順序付き複素波配列 \\
\(\mathcal{W}\) & 波の置換で同一視した無名複素波多重集合 \\
Q & 複素二乗和による閉塞量 \\
q & 有限回帰次数 \\
s & 数値計算上の更新順序。物理時間ではない \\
F & 無名集合から抽出された因子 \\
C & 因子を構成する部分多重集合 \\
\(\theta\) & 因子の潜在パラメータ \\
\(\varepsilon\) & 因子再構成残差 \\
\(\Omega_{\mathrm{fact}}\) & 許容対称性を除いた非同値因数分解数 \\
D & シンプレックス再構成に用いる無名複素二乗量集合 \\
n & 再構成候補となるシンプレックス頂点数 \\
G & 候補割当てから構成した複素Gram行列 \\
k & 倍音因子の整数次数 \\
\(\phi_0\) & 復元された共通位相中心 \\
\(\Sigma\) & 因子の不変署名 \\
\end{longtable}
}

\begin{center}\rule{0.5\linewidth}{0.5pt}\end{center}

\subsection{付録B：事前登録すべき実験項目}\label{ux4ed8ux9332bux4e8bux524dux767bux9332ux3059ux3079ux304dux5b9fux9a13ux9805ux76ee}

各数値実験を実施する前に、少なくとも次を固定する。

\begin{enumerate}
\def\labelenumi{\arabic{enumi}.}
\tightlist
\item
  入力として許可するデータ項目
\item
  復元器から秘匿する生成情報
\item
  同一視する対称変換
\item
  因子残差の定義
\item
  数値許容値
\item
  探索範囲と打切り条件
\item
  一意性の判定方法
\item
  正解回収の評価方法
\item
  ランダム対照群の生成方法
\item
  摂動分布
\item
  成功条件
\item
  反証条件
\end{enumerate}

この事前登録により、結果を見た後で因子定義、許容値または探索範囲を都合よく変更する循環を防ぐ。

\end{document}
